\documentclass[UTF8, 12pt, fontset=fandol, oneside]{ctexart}
\usepackage{researchnotes}

% 保留分册独立编译时的章号前缀。
\renewcommand{\thesection}{13.\arabic{section}}

\title{第十三章\quad 多元函数的极限和连续}
\author{}
\date{}

\begin{document}

\maketitle

在本套教材的第一册与第二册中, 我们已经系统地学习了一元微积分与级数理论. 但在理论与实践中, 仅仅一元函数远远不能满足需要. 这是因为, 在许多事物的变化过程中, 一个变量的变化过程往往依赖于多个变量. 就拿我们每天生活的空间来说, 它是一个三维的立体空间, 因此几乎所有跟空间位置有关的变量一般都要用空间点的坐标来描述, 从而它们就不太可能用一元函数来刻画。

另外, 即使在数学研究中, 由于一元函数的研究仅仅是局限于数轴 $\mathbb R$ 的子集上定义的函数, 它们基本上已不再是现代数学研究的主要对象. 在当今的数学研究中, 大部分的研究对象都是关于高维空间 ($n(n\ge2)$ 维空间) 的一些问题. 因此, 我们对多元函数 (映射) 的学习是十分必要的。

多元微积分的主要内容是将一元函数的微积分理论推广到高维空间上的多元函数. 大家会发现, 我们将平行于一元微积分的基本理论来研究多元微积分. 值得指出的是, 由于多元函数的微积分理论是建立在一元微积分的基础之上的, 读者如果具备一元微积分的坚实基础, 并且有较好的空间想象能力, 就能学好多元微积分。

\section{欧氏空间 $\mathbb R^n$}

\subsection{欧氏空间 $\mathbb R^n$}

在本节中, 我们先来讨论多元函数定义域的问题. 多元函数的定义域是高维空间的子集, 这些子集相对于 $\mathbb R$ 中的子集将更为复杂. 为了研究高维空间的子集, 我们必须先研究一下它们所在的空间
\[
  \mathbb R^n=\{(x_1,x_2,\cdots,x_n):x_i\in\mathbb R,\ i=1,2,\cdots,n\}
  =\underbrace{\mathbb R\times\mathbb R\times\cdots\times\mathbb R}_{n\text{ 个}}.
\]

读者应特别注意当 $n=2$ 时的情形. 多元微积分与一元微积分的许多本质区别将在 $n=1$ 和 $n=2$ 两者之间发生, 对于 $n>2$ 的情形则与 $n=2$ 的情形具有很大的相似性。

在以下讨论中我们将假定 $n\ge2$. 记 $x=(x_1,x_2,\cdots,x_n)\in\mathbb R^n$, 我们称 $x$ 为 $\mathbb R^n$ 中的一个点或向量, $x_i(i=1,2,\cdots,n)$ 称为 $x$ 的第 $i$ 个坐标或分量. 今后, 我们也常常用列向量 $(x_1,x_2,\cdots,x_n)^{\mathrm T}$ 来表示同一个 $x$, 这里 $(x_1,x_2,\cdots,x_n)^{\mathrm T}$ 是 $(x_1,x_2,\cdots,x_n)$ 的转置. 另外, 记 $\mathbf 0=(0,0,\cdots,0)$ 为 $\mathbb R^n$ 中的原点或零向量. 在代数课程中, 我们已经在 $\mathbb R^n$ 中引进了加法与数乘运算. 由于在今后要经常用到它们, 在这里我们做一下简单介绍。

$\mathbb R^n$ 中的加法运算定义如下: 设 $x=(x_1,x_2,\cdots,x_n), y=(y_1,y_2,\cdots,y_n)\in\mathbb R^n$, 定义
\[
  x+y=(x_1+y_1,x_2+y_2,\cdots,x_n+y_n)\in\mathbb R^n,
\]
并称 $x+y$ 为 $x$ 与 $y$ 的和。

$\mathbb R^n$ 中的数乘运算定义如下: 设 $\alpha\in\mathbb R, x=(x_1,x_2,\cdots,x_n)\in\mathbb R^n$, 定义
\[
  \alpha x=(\alpha x_1,\alpha x_2,\cdots,\alpha x_n)\in\mathbb R^n,
\]
并称 $\alpha x$ 为 $\alpha$ 与 $x$ 的数乘。

这两种运算称为 $\mathbb R^n$ 中的线性运算. 在 $n=2,3$ 时, 它们具有鲜明的几何意义. 对于 $\forall x,y,z\in\mathbb R^n,\alpha,\beta\in\mathbb R$, 容易验证它们满足:
\[
\begin{gathered}
  (1)\quad \text{\textbf{交换律}}\quad x+y=y+x;\\
  (2)\quad \text{\textbf{结合律}}\quad (x+y)+z=x+(y+z),\quad (\alpha\beta)x=\alpha(\beta x);\\
  (3)\quad \text{\textbf{分配律}}\quad \alpha(x+y)=\alpha x+\alpha y,\quad
  (\alpha+\beta)x=\alpha x+\beta x.
\end{gathered}
\]
另外, 在加法运算中, 存在零元素 $\mathbf 0\in\mathbb R^n$, 它满足: 对于 $\forall x\in\mathbb R^n$, 有
\[
  x+\mathbf 0=x;
\]
在数乘运算中, 存在单位元 $1\in\mathbb R$, 它满足: 对于 $\forall x\in\mathbb R^n$, 有
\[
  1x=x.
\]

对 $\mathbb R^n$ 赋予上述线性运算后, 我们称 $\mathbb R^n$ 为一个 $n$ 维向量空间 (简称空间). 在这个空间中, 它还有一个重要的运算——内积运算, 它的定义如下:

设 $x=(x_1,x_2,\cdots,x_n), y=(y_1,y_2,\cdots,y_n)\in\mathbb R^n$, 则 $x$ 与 $y$ 的内积定义为
\[
  xy=\sum_{i=1}^n x_iy_i\in\mathbb R.
\]
从内积的定义容易看出它具有以下一些基本性质:
\[
\begin{aligned}
  &(1)\quad \text{\textbf{正定性}}\quad \text{对于 } \forall x\in\mathbb R^n,\ \text{有 } xx\ge0,\ \text{并且上述等号成立当且仅当 } x=\mathbf 0;\\
  &(2)\quad \text{\textbf{对称性}}\quad \text{对于 } \forall x,y\in\mathbb R^n,\ \text{有 } xy=yx;\\
  &(3)\quad \text{对于 } \forall x,y,z\in\mathbb R^n,\ \text{有 } x(y+z)=xy+xz;\\
  &(4)\quad \text{对于 } \forall\alpha\in\mathbb R \text{ 和 } \forall x,y\in\mathbb R^n,\ \text{有 }(\alpha x)y=\alpha(xy).
\end{aligned}
\]
向量空间 $\mathbb R^n$ 有了内积运算后, 我们称 $\mathbb R^n$ 为欧几里得 (Euclid) 空间或欧氏空间. 利用内积运算, 我们定义向量 $x\in\mathbb R^n$ 的模如下:
\[
  \|x\|=\sqrt{xx}=\sqrt{\sum_{i=1}^n x_i^2}.
\]
在代数学中, 我们已经知道, $\mathbb R^n(n\ge2)$ 中两个非零向量 $x$ 与 $y$ 的内积
\[
  xy=\|x\|\cdot\|y\|\cos(x,y),
\]
其中 $(x,y)$ 是向量 $x$ 与 $y$ 的夹角. 由此我们可以清楚地知道内积的几何意义。

利用向量的模, 我们可以给出 $\mathbb R^n$ 中两个点之间的距离的定义。

\noindent\textbf{定义 13.1.1}\quad 设 $x=(x_1,x_2,\cdots,x_n)$ 与 $y=(y_1,y_2,\cdots,y_n)$ 为 $\mathbb R^n$ 中任意两个点, 则 $x$ 与 $y$ 的距离定义为
\[
  |x-y|=\|x-y\|=\sqrt{\sum_{i=1}^n(x_i-y_i)^2}.
\]
显然, 在 $\mathbb R,\mathbb R^2$ 或 $\mathbb R^3$ 中, 两个点 $x$ 与 $y$ 的距离即是连接 $x$ 与 $y$ 的线段的长度. 请读者注意, 我们用 $|x-y|$ (而不是用 $\|x-y\|$) 来表示 $x,y$ 之间的距离主要是为了记号简便, 且这个记号与 $\mathbb R$ 中两点之间距离的相应记号保持一致. 由于对于 $\forall x\in\mathbb R^n$, 我们有 $\|x\|=|x-\mathbf 0|=|x|$, 因此, 今后我们也用 $|x|$ 来记 $x$ 的模。

从距离的定义容易推出它满足以下的性质:
\[
\begin{aligned}
  &(1)\quad \text{\textbf{正定性}}\quad
  \text{对于 } \forall x,y\in\mathbb R^n,\ \text{有 } |x-y|\ge0,\ \text{并且 } |x-y|=0 \text{ 的充分必要条件是 } x=y;\\
  &(2)\quad \text{\textbf{对称性}}\quad
  \text{对于 } \forall x,y\in\mathbb R^n,\ \text{有 } |x-y|=|y-x|;\\
  &(3)\quad \text{\textbf{三角不等式}}\quad
  \text{对于 } \forall x,y,z\in\mathbb R^n,\ \text{有 } |x-z|\le |x-y|+|y-z|.
\end{aligned}
\]
距离与向量的模两个概念是可以相互转化的, 因此我们也称定义了距离后的空间 $\mathbb R^n$ 为欧氏空间. 当 $n=2$ 时, 我们常用 $(x,y)$ 来表示平面 $\mathbb R^2$ 中的点; 当 $n=3$ 时, 用 $(x,y,z)$ 来表示空间 $\mathbb R^3$ 中的点. 今后在 $\mathbb R^2$ 中, 为了记号方便, 我们也用 $\boldsymbol{i}$ 来表示单位向量 $(1,0)$, 用 $\boldsymbol{j}$ 来表示单位向量 $(0,1)$, 并称它们为单位坐标向量. 因此, 对于 $\forall(x,y)\in\mathbb R^2$, 我们有
\[
  (x,y)=x\boldsymbol{i}+y\boldsymbol{j}.
\]
相应地, 在 $\mathbb R^3$ 中, 我们则记 $\boldsymbol{i}=(1,0,0), \boldsymbol{j}=(0,1,0), \boldsymbol{k}=(0,0,1)$. 于是对于 $\forall(x,y,z)\in\mathbb R^3$, 有
\[
  (x,y,z)=x\boldsymbol{i}+y\boldsymbol{j}+z\boldsymbol{k}.
\]

\noindent\textbf{例 13.1.1}\quad 设
\[
  A=
  \begin{pmatrix}
    a_{11} & a_{12} & \cdots & a_{1n}\\
    a_{21} & a_{22} & \cdots & a_{2n}\\
    \vdots & \vdots &        & \vdots\\
    a_{m1} & a_{m2} & \cdots & a_{mn}
  \end{pmatrix}
\]
是一个 $m\times n$ 矩阵, 其中 $a_{ji}\in\mathbb R(j=1,2,\cdots,m;\ i=1,2,\cdots,n)$. 定义
\[
  \|A\|=\left(\sum_{j=1}^m\sum_{i=1}^n a_{ji}^2\right)^{\frac12}
\]
对于 $\forall x=(x_1,x_2,\cdots,x_n)^{\mathrm T}\in\mathbb R^n$, 证明 $|Ax|\le \|A\||x|$。

\noindent\textbf{证明}\quad 利用柯西-施瓦茨不等式我们有
\[
\begin{aligned}
  |Ax|^2
  &=\sum_{j=1}^m\left(\sum_{i=1}^n a_{ji}x_i\right)^2
    \le \sum_{j=1}^m\left(\sum_{i=1}^n a_{ji}^2\sum_{i=1}^n x_i^2\right)\\
  &=\left(\sum_{j=1}^m\sum_{i=1}^n a_{ji}^2\right)
    \left(\sum_{i=1}^n x_i^2\right)
    =(\|A\||x|)^2.
\end{aligned}
\]
对上面不等式的两边分别开方即得所证。

\noindent\textbf{注}\quad $\|A\|=\left(\sum\limits_{j=1}^m\sum\limits_{i=1}^n a_{ji}^2\right)^{\frac12}$ 称为矩阵 $A$ 的范数, 在本书后面章节我们还会遇到它。

\subsection{点列极限}

下面我们给出欧氏空间 $\mathbb R^n$ 中邻域的概念。

\noindent\textbf{定义 13.1.2}\quad 设 $x_0=(x_1^0,x_2^0,\cdots,x_n^0)\in\mathbb R^n,\delta>0$, 称集合
\[
  U(x_0,\delta)=\{x=(x_1,x_2,\cdots,x_n)\in\mathbb R^n:|x-x_0|<\delta\}
\]
为以 $x_0$ 为心的 $\delta$ 邻域; 称集合 $U_0(x_0,\delta)=U(x_0,\delta)\setminus\{x_0\}$ 为 $x_0$ 的 $\delta$ 去心邻域。

上述定义的邻域通常称为球形邻域. 我们经常要用到的另外一种邻域是方形邻域. 设 $x_0=(x_1^0,x_2^0,\cdots,x_n^0)\in\mathbb R^n,\delta>0$, 定义
\[
  N(x_0,\delta)=\{x=(x_1,x_2,\cdots,x_n): |x_i-x_i^0|<\delta,\ i=1,2,\cdots,n\};
\]
并称它为 $x_0$ 的方形邻域. 容易证明这两种邻域有下面的包含关系:
\[
  U(x_0,\delta)\subset N(x_0,\delta)\subset U(x_0,\sqrt n\delta).
\]
另外, 我们称 $N_0(x_0,\delta)=N(x_0,\delta)\setminus\{x_0\}$ 为方形去心邻域。

以上所定义的邻域中, 当 $n=1$ 时, 球形邻域与方形邻域是一致的, 它们就是实数轴上的开区间. 当 $n=2$ 时, 球形邻域就是开圆盘, 方形邻域就是开正方形. 当 $n=3$ 时, 球形邻域就是开球, 方形邻域就是开立方体。

下面我们利用距离来描述 $\mathbb R^n$ 中点列的极限. 如果点列 $\{x_k\}$ 中的点 $x_k$ 到某个点 $x_0$ 的距离趋于零, 即
\[
  \lim_{k\to\infty}|x_k-x_0|=0,
\]
那么我们就说点列 $\{x_k\}$ 收敛于 $x_0$. 用邻域的语言可以给出如下定义。

\noindent\textbf{定义 13.1.3}\quad 设 $\{x_k\}$ 是 $\mathbb R^n$ 中的一个点列, 若存在 $x_0\in\mathbb R^n$, 使得对于 $\forall\varepsilon>0,\exists K\in\mathbb N$, 当 $k>K$ 时, 有 $|x_k-x_0|<\varepsilon$, 即 $x_k\in U(x_0,\varepsilon)$, 则称 $\{x_k\}$ 是收敛点列, 并称 $\{x_k\}$ 收敛于 $x_0$, 记做
\[
  \lim_{k\to\infty}x_k=x_0.
\]
这时也称 $x_0$ 为 $\{x_k\}$ 的极限. 若不存在 $x_0\in\mathbb R^n$, 使得 $\lim\limits_{k\to\infty}|x_k-x_0|=0$, 则称 $\{x_k\}$ 发散。

在本章中, 为了避免与 $\mathbb R^n$ 中的 $n$ 混淆, 我们常用 $\{x_k\}$ 来表示 $\mathbb R^n$ 中的点列, 而不用 $\{x_n\}$. 设 $x_0=(x_1^0,x_2^0,\cdots,x_n^0)$, $x_k=(x_1^k,x_2^k,\cdots,x_n^k)$, 则有下面的定理。

\noindent\textbf{定理 13.1.1}\quad 设 $\{x_k\}$ 是 $\mathbb R^n$ 中的一个点列, $x_0=(x_1^0,x_2^0,\cdots,x_n^0)\in\mathbb R^n$, 则 $\lim\limits_{k\to\infty}x_k=x_0$ 的充分必要条件是: 对于 $\forall i(1\le i\le n)$, 有
\[
  \lim_{k\to\infty}x_i^k=x_i^0.
\]

\noindent\textbf{证明}\quad 对于 $\forall i=1,2,\cdots,n$ 和 $\forall k\in\mathbb N$, 有
\[
  |x_i^k-x_i^0|\le |x_k-x_0|\le \sum_{j=1}^n|x_j^k-x_j^0|.
\]
由此不等式容易推出定理 13.1.1. 证毕。

若 $\mathbb R^n$ 中的集合 $E$ 包含于某个球形邻域 $U(0,M)$ 中, 则称 $E$ 是有界集. 若点列 $\{x_k\}$ 的所有点组成的集合是有界集, 则称 $\{x_k\}$ 是有界点列. 与数列的情形类似, $\mathbb R^n$ 中的收敛点列具有下面的性质。

\noindent\textbf{性质 13.1.2}\quad 收敛点列的极限是唯一的。

\noindent\textbf{性质 13.1.3}\quad 收敛点列一定是有界的。

\noindent\textbf{性质 13.1.4}\quad 若 $\mathbb R^n$ 中的点列 $\{x_k\},\{y_k\}$ 满足
\[
  \lim_{k\to\infty}x_k=x_0,\qquad \lim_{k\to\infty}y_k=y_0,
\]
且 $\alpha,\beta\in\mathbb R$, 则
\[
  (1)\quad \lim_{k\to\infty}(\alpha x_k+\beta y_k)=\alpha x_0+\beta y_0;
\]
\[
  (2)\quad \lim_{k\to\infty}x_ky_k=x_0y_0.
\]
对于以上性质, 我们只证性质 13.1.4 的 (2), 读者应该注意这个性质说的是两个收敛点列的内积收敛到点列极限的内积。

\noindent\textbf{证明}\quad 设 $x_k=(x_1^k,x_2^k,\cdots,x_n^k)$, $y_k=(y_1^k,y_2^k,\cdots,y_n^k)$, $x_0=(x_1^0,x_2^0,\cdots,x_n^0)$, $y_0=(y_1^0,y_2^0,\cdots,y_n^0)$. 由定理 13.1.1 知, 对于 $\forall i=1,2,\cdots,n$, 有
\[
  \lim_{k\to\infty}x_i^k=x_i^0,\qquad \lim_{k\to\infty}y_i^k=y_i^0.
\]
由数列极限的性质知
\[
  \lim_{k\to\infty}x_i^ky_i^k=x_i^0y_i^0.
\]
于是
\[
  \lim_{k\to\infty}x_ky_k
  =\lim_{k\to\infty}\sum_{i=1}^n x_i^ky_i^k
  =\sum_{i=1}^n x_i^0y_i^0=x_0y_0.
\]
证毕。

性质 13.1.4 的 (2) 是实数乘积极限性质的推广, 但当 $n\ge2$ 时, 这个性质与实数乘积极限性质有明显的差别. 在 $\mathbb R^n$ 中, 两个非零向量的内积可能等于零, 例如两个非零向量正交时。

\noindent\textbf{例 13.1.2}\quad 设
\[
  x_k=\left(\frac{\sqrt k\sin k}{k},\,k^4e^{-k^2},\,\cos k\ln\left(1+\frac1k\right),\,k\tan\frac1k\right)\in\mathbb R^4
\]
$(k=1,2,\cdots)$, 求 $\lim\limits_{k\to\infty}x_k$。

\noindent\textbf{解}\quad 由于
\[
  \lim_{k\to\infty}\frac{\sqrt k\sin k}{k}=0,\qquad
  \lim_{k\to\infty}k^4e^{-k^2}=0,
\]
\[
  \lim_{k\to\infty}\cos k\ln\left(1+\frac1k\right)=0,\qquad
  \lim_{k\to\infty}k\tan\frac1k=1,
\]
故
\[
  \lim_{k\to\infty}x_k=(0,0,0,1).
\]

\subsection{聚点}

聚点是极限理论中重要的概念. 我们曾在 $\mathbb R$ 中讨论过它(见第一册). 对于 $\mathbb R^n$ 的一般情形, 聚点的定义如下:

\noindent\textbf{定义 13.1.4}\quad 设 $E\subset\mathbb R^n$. 若 $x\in\mathbb R^n$ 的每个 $\delta$ 邻域 $U(x,\delta)(\delta>0)$ 中都含有异于 $x$ 的 $E$ 中的点, 则称 $x$ 是 $E$ 的聚点或极限点。

在上述定义中, 我们也可以用方形邻域来代替球形邻域. 容易看出, 若 $x$ 是 $E$ 的聚点, 则 $x$ 的每个邻域中都含有无穷多个 $E$ 中的点. 另外, $x$ 可以属于 $E$, 也可以不属于 $E$。

若 $x$ 不是 $E$ 的聚点, 则存在 $\delta_0>0$, 使得 $U_0(x,\delta_0)$ 中不含有 $E$ 中的点. 如果 $x\in E$ 且 $x$ 不是 $E$ 的聚点, 则称 $x$ 是 $E$ 的孤立点. 这时存在 $\delta_0>0$, 使得
\[
  U(x,\delta_0)\cap E=\{x\}.
\]

利用点列的极限, 我们可以描述一个集合的聚点。

\noindent\textbf{定理 13.1.5}\quad 设 $\varnothing\ne E\subset\mathbb R^n$, 则 $x$ 是 $E$ 的聚点的充分必要条件是: 存在 $E$ 中的互异点列 $\{x_k\}$, 使得
\[
  \lim_{k\to\infty}x_k=x.
\]

定理 13.1.5 的证明与 $\mathbb R$ 中的情形类似, 请读者自己给出。

\noindent\textbf{例 13.1.3}\quad 设
\[
  E=\{(x_1,x_2,\cdots,x_n)\in\mathbb R^n:\text{对于 }\forall i(1\le i\le n),\ x_i\text{ 是有理数}\}.
\]
求 $E$ 的聚点集。

\noindent\textbf{解}\quad 设 $x_0=(x_1^0,x_2^0,\cdots,x_n^0)\in\mathbb R^n$. 对于 $\forall\delta>0$, 考虑 $x_0$ 的方形邻域 $N(x_0,\delta)$. 对于每一个 $i$, 开区间 $(x_i^0-\delta,x_i^0+\delta)$ 中都含有有理数 $x_i'$ 且 $x_i'\ne x_i^0$. 于是
\[
  (x_1',x_2',\cdots,x_n')\in N(x_0,\delta)\cap E,
\]
且 $(x_1',x_2',\cdots,x_n')\ne x_0$. 这说明 $x_0$ 是 $E$ 的聚点. 由于 $x_0$ 是 $\mathbb R^n$ 中任意一点, 故 $E$ 的聚点集是 $\mathbb R^n$。

如果集合 $A$ 中的每个点的任意邻域中都含有集合 $B$ 中的点, 则称 $B$ 在 $A$ 中稠密. 上面的例子表明, $E$ 在 $\mathbb R^n$ 中稠密. 请读者利用例 13.1.3 的结论, 将 $E$ 中的点排成一个点列。

\subsection{开集与闭集}

设 $E\subset\mathbb R^n$, 记 $E$ 的余集为 $E^c=\mathbb R^n\setminus E$. 下面我们利用邻域来对 $\mathbb R^n$ 中的点进行分类。

\noindent\textbf{定义 13.1.5}\quad 设 $E\subset\mathbb R^n,\ x\in\mathbb R^n$。

(1) 若存在 $\delta>0$, 使得 $U(x,\delta)\subset E$, 则称 $x$ 是 $E$ 的内点. $E$ 的全体内点组成的集合记为 $E^\circ$, 称为 $E$ 的内部。

(2) 若存在 $\delta>0$, 使得 $U(x,\delta)\cap E=\varnothing$, 则称 $x$ 是 $E$ 的外点. $E$ 的全体外点组成的集合称为 $E$ 的外部。

(3) 若对于 $\forall\delta>0$, 都有 $U(x,\delta)\cap E\ne\varnothing$ 且 $U(x,\delta)\cap E^c\ne\varnothing$, 则称 $x$ 是 $E$ 的边界点. $E$ 的全体边界点组成的集合记为 $\partial E$, 称为 $E$ 的边界。

由定义可知, $E$ 的内点一定属于 $E$; $E$ 的外点一定不属于 $E$, 且 $E$ 的外部就是 $(E^c)^\circ$; $E$ 的边界点可能属于 $E$, 也可能不属于 $E$. 一个点 $x$ 是 $E$ 的边界点的充分必要条件是: $x$ 既不是 $E$ 的内点, 也不是 $E$ 的外点。

\noindent\textbf{例 13.1.4}\quad 设
\[
  E=\{(x,y):x^2+y^2<1,\ y\ge x\},
\]
求 $E^\circ,(E^c)^\circ$ 和 $\partial E$。

\noindent\textbf{解}\quad 参见图 13.1.1. 容易看出
\[
  E^\circ=\{(x,y):x^2+y^2<1,\ y>x\}=E\setminus\{(x,y):y=x\},
\]
\[
  (E^c)^\circ=\mathbb R^2\setminus\{(x,y):x^2+y^2\le1,\ y\ge x\},
\]
\[
  \partial E=\{(x,y):x^2+y^2=1,\ y\ge x\}\cup\{(x,y):x^2+y^2<1,\ y=x\}.
\]

在一元微积分中, 开区间与闭区间起着重要的作用. 在 $\mathbb R^n$ 中与它们密切相关的概念是开集与闭集. 开集与闭集是 $\mathbb R^n$ 中两类重要的子集, 它们在多元函数的极限与连续理论中有重要作用。

\noindent\textbf{定义 13.1.6}\quad 设 $E\subset\mathbb R^n$. 若 $E=E^\circ$, 则称 $E$ 是开集. 规定空集 $\varnothing$ 是开集。

显然, 每个球形邻域和方形邻域都是开集. 在 $\mathbb R$ 中, 开集可以表示为至多可数个互不相交的开区间的并. 在 $\mathbb R^2$ 中, 开集可以看成是不含边界的平面点集. 例如
\[
  E=\{(x,y): |(x,y)-(0,k)|<1/2,\ k\in\mathbb Z\}
\]
是一个开集, 它由无穷多个两两不相交的开圆盘组成. 即使一个开集只有一块, 它的边界也可能很复杂. 例如
\[
  E=\{(x,y):0\le y\le 1/2,\ x=1/k,\ k=1,2,\cdots\},
\]
则
\[
  D=((0,1)\times(0,1))\setminus E
\]
的边界为 $E\cup\partial([0,1]\times[0,1])$, 且 $D$ 是开集。

\noindent\textbf{定理 13.1.6}\quad $\mathbb R^n$ 中的开集具有以下性质:

(1) $\mathbb R^n$ 和 $\varnothing$ 都是开集;

(2) 任意多个开集的并是开集;

(3) 有限多个开集的交是开集。

定理 13.1.6 的证明很容易, 请读者自己完成。

需要注意的是, 任意多个开集的交不一定是开集. 事实上, 对于 $x\in\mathbb R^n$, 有
\[
  \{x\}=\bigcap_{k=1}^{+\infty}U\left(x,\frac1k\right).
\]
在 $\mathbb R$ 中, 闭区间的余集是开集. 由此我们可以引入闭集的概念。

\noindent\textbf{定义 13.1.7}\quad 设 $E\subset\mathbb R^n$. 若 $E^c$ 是开集, 则称 $E$ 是闭集。

在 $\mathbb R$ 中, 闭集不一定能表示为至多可数个闭区间的并. 由定义可知, 对于任意开集 $E$, 集合 $F=E\cup\partial E$ 是闭集。

\noindent\textbf{定理 13.1.7}\quad $\mathbb R^n$ 中的闭集具有以下性质:

(1) $\mathbb R^n$ 和 $\varnothing$ 都是闭集;

(2) 有限多个闭集的并是闭集;

(3) 任意多个闭集的交是闭集。

定理 13.1.7 可以直接由闭集的定义和定理 13.1.6 得到. 事实上, 对任意一族集合 $\{E_\lambda\}_{\lambda\in\Lambda}$, 我们有德摩根公式
\[
  \left(\bigcup_{\lambda\in\Lambda}E_\lambda\right)^c
  =\bigcap_{\lambda\in\Lambda}E_\lambda^c,\qquad
  \left(\bigcap_{\lambda\in\Lambda}E_\lambda\right)^c
  =\bigcup_{\lambda\in\Lambda}E_\lambda^c.
\]

设 $E\subset\mathbb R^n$, 记 $E$ 的全体聚点组成的集合为 $E'$, 称为 $E$ 的导集. 集合
\[
  \overline E=E\cup E'
\]
称为 $E$ 的闭包。

\noindent\textbf{定理 13.1.8}\quad 设 $E\subset\mathbb R^n$, 则 $E$ 是闭集的充分必要条件是 $E=\overline E$。

\noindent\textbf{证明}\quad 由闭包的定义可知, 只需证明 $E$ 是闭集的充分必要条件是 $E'\subset E$。

充分性. 若 $E$ 的所有聚点都属于 $E$, 则 $E^c$ 中没有 $E$ 的聚点. 于是对于 $\forall x\in E^c$, 存在 $\delta_x>0$, 使得
\[
  U(x,\delta_x)\cap E=\varnothing.
\]
这说明 $U(x,\delta_x)\subset E^c$, 因而 $E^c$ 是开集, 即 $E$ 是闭集。

必要性. 若 $E$ 是闭集, 则 $E^c$ 是开集. 对于 $\forall x\in E^c$, 存在 $\delta_x>0$, 使得 $U(x,\delta_x)\subset E^c$, 因而 $x$ 不是 $E$ 的聚点. 这说明 $E^c$ 中没有 $E$ 的聚点, 即 $E'\subset E$, 从而 $E=E\cup E'=\overline E$. 证毕。

\subsection{欧氏空间 $\mathbb R^n$ 中的基本定理}

本小节我们将实直线中的一些基本定理推广到欧氏空间 $\mathbb R^n$ 中, 包括完备性定理、闭集套定理、聚点原理和有限覆盖定理。

\noindent\textbf{1. 完备性}

设 $\{x_k\}$ 是 $\mathbb R^n$ 中的点列. 若对于 $\forall\varepsilon>0,\exists K\in\mathbb N$, 当 $k',k''>K$ 时, 有
\[
  |x_{k'}-x_{k''}|<\varepsilon,
\]
则称 $\{x_k\}$ 是柯西点列. 若空间 $X\subset\mathbb R^n$ 中的每个柯西点列都收敛, 且极限属于 $X$, 则称 $X$ 是完备的。

\noindent\textbf{定理 13.1.9}\quad 欧氏空间 $\mathbb R^n$ 是完备的。

\noindent\textbf{证明}\quad 设 $x_k=(x_1^k,x_2^k,\cdots,x_n^k)\in\mathbb R^n(k=1,2,\cdots)$ 是一个柯西点列. 则对于 $\forall i=1,2,\cdots,n$, 数列 $\{x_i^k\}$ 是 $\mathbb R$ 中的柯西数列. 由 $\mathbb R$ 的完备性知, 每个数列 $\{x_i^k\}$ 都收敛. 由定理 13.1.1 可知 $\{x_k\}$ 收敛. 证毕。

完备性是欧氏空间 $\mathbb R^n$ 的一个重要性质, 在数学分析中有广泛应用。

\noindent\textbf{2. 闭集套定理}

设 $E$ 是 $\mathbb R^n$ 中的有界集, 定义
\[
  \operatorname{diam}(E)=\sup_{x,y\in E}\{|x-y|\},
\]
称为 $E$ 的直径。

\noindent\textbf{定理 13.1.10}\quad 设 $F_k\subset\mathbb R^n(k=1,2,\cdots)$ 是一列非空闭集, 且满足

(1) 对于 $\forall k\in\mathbb N$, 有 $F_k\supset F_{k+1}$;

(2) $\lim\limits_{k\to\infty}\operatorname{diam}(F_k)=0$。

则存在唯一的 $x_0\in\mathbb R^n$, 使得 $x_0\in F_k(k=1,2,\cdots)$, 即
\[
  \{x_0\}=\bigcap_{k=1}^{+\infty}F_k.
\]

\noindent\textbf{证明}\quad 对于每个 $k$, 任取 $x_k\in F_k$. 下面证明 $\{x_k\}$ 是柯西点列. 对于 $\forall\varepsilon>0$, 存在 $K\in\mathbb N$, 当 $k>K$ 时, 有 $\operatorname{diam}(F_k)<\varepsilon$. 若 $k''>k'>K$, 则 $x_{k'}\in F_{k'}$, $x_{k''}\in F_{k''}\subset F_{k'}$, 因而
\[
  |x_{k'}-x_{k''}|\le \operatorname{diam}(F_{k'})<\varepsilon.
\]
这说明 $\{x_k\}$ 是柯西点列. 由定理 13.1.9 知 $\{x_k\}$ 收敛, 记 $x_0=\lim\limits_{k\to\infty}x_k$. 对于任意给定的 $k_0$, 当 $k>k_0$ 时, 有 $x_k\in F_{k_0}$. 由于 $F_{k_0}$ 是闭集, 故 $x_0\in F_{k_0}$. 由 $k_0$ 的任意性知 $x_0\in\bigcap\limits_{k=1}^{+\infty}F_k$。

若另有 $x'\in\bigcap\limits_{k=1}^{+\infty}F_k$, 则
\[
  |x'-x_0|\le \lim_{k\to\infty}\operatorname{diam}(F_k)=0.
\]
于是 $x'=x_0$. 证毕。

实直线上的闭区间套定理是定理 13.1.10 的特殊情形. 定理 13.1.10 中的闭集不能换成开集; 如果 $F_k$ 不是闭集, 结论可能不成立。

\noindent\textbf{3. 聚点原理}

在 $\mathbb R^n$ 中, 我们同样有以下的聚点原理。

\noindent\textbf{定理 13.1.11}\quad $\mathbb R^n$ 中的有界点列必有收敛子列。

\noindent\textbf{证明}\quad 设 $\{x_k\}$ 是 $\mathbb R^n$ 中的有界点列, 其中 $x_k=(x_1^k,x_2^k,\cdots,x_n^k)$. 则对于每个 $i=1,2,\cdots,n$, 数列 $\{x_i^k\}$ 都有界. 由数列的聚点原理知, $\{x_1^k\}$ 有收敛子列 $\{x_1^{k_j'}\}$. 再从 $\{x_2^{k_j'}\}$ 中取收敛子列 $\{x_2^{k_j''}\}$. 依此进行, 经过 $n$ 次后得到 $\{x_k\}$ 的一个子列 $\{x_{k_j}\}$, 使得它的每个分量数列都收敛. 因而由定理 13.1.1 知 $\{x_{k_j}\}$ 收敛. 证毕。

\noindent\textbf{定理 13.1.12}\quad $\mathbb R^n$ 中的任一有界无限集 $E$ 至少有一个聚点。

\noindent\textbf{证明}\quad 由于 $E$ 是无限集, 可从 $E$ 中取出互异点列 $\{x_k\}$. 又由于 $E$ 有界, 故 $\{x_k\}$ 是有界点列. 由定理 13.1.11 知 $\{x_k\}$ 有收敛子列. 设该收敛子列的极限为 $x_0$. 由于 $\{x_k\}$ 中的点互异, 故 $x_0$ 是集合 $\{x_k\}$ 的聚点, 从而也是 $E$ 的聚点. 证毕。

\noindent\textbf{4. 紧集}

在实直线上, 有限覆盖定理是一个重要定理. 在 $\mathbb R^n$ 中, 与有限覆盖定理密切相关的概念是紧集。

设 $E\subset\mathbb R^n$, $\{O_\lambda\}_{\lambda\in\Lambda}$ 是 $\mathbb R^n$ 中的一族开集. 若
\[
  E\subset \bigcup_{\lambda\in\Lambda}O_\lambda,
\]
则称 $\{O_\lambda\}_{\lambda\in\Lambda}$ 是 $E$ 的一个开覆盖. 若 $\Lambda$ 是有限集, 则称它为有限开覆盖。

\noindent\textbf{定义 13.1.8}\quad 设 $E\subset\mathbb R^n$. 若 $E$ 的每一个开覆盖都存在有限子覆盖, 即对于 $E$ 的任意开覆盖 $\{O_\lambda\}_{\lambda\in\Lambda}$, 都存在其中有限个开集 $O_1,O_2,\cdots,O_K(K<\infty)$, 使得
\[
  E\subset\bigcup_{k=1}^K O_k,
\]
则称 $E$ 是紧集。

紧集的定义强调的是: $E$ 的每一个开覆盖都存在有限子覆盖, 而不是只要找到某一个有限开覆盖即可。

在 $\mathbb R$ 中, 有限个有界闭区间的并是紧集. 若 $E$ 是 $\mathbb R^n$ 中的无界集, 则 $E$ 不是紧集. 事实上, $\{U(0,k)\}_{k\in\mathbb N}$ 是 $E$ 的一个开覆盖, 但它不存在有限子覆盖. 例如 $E=\{1/k:k\in\mathbb N\}$ 不是紧集; 但 $E\cup E'=\{1/k:k\in\mathbb N\}\cup\{0\}$ 是紧集。

\noindent\textbf{定理 13.1.13}\quad 设 $E\subset\mathbb R^n$, 则 $E$ 是紧集的充分必要条件是 $E$ 是有界闭集。

\noindent\textbf{证明}\quad 必要性. 对于任意集合 $E\subset\mathbb R^n$, $\{U(0,k)\}_{k\in\mathbb N}$ 是 $E$ 的一个开覆盖. 由于 $E$ 是紧集, 故它存在有限子覆盖 $U(0,k_1),U(0,k_2),\cdots,U(0,k_J)$, 不妨设 $0<k_1<k_2<\cdots<k_J<\infty$. 于是
\[
  E\subset\bigcup_{j=1}^J U(0,k_j)=U(0,k_J),
\]
故 $E$ 有界。

下面证明 $E$ 是闭集. 反设 $E$ 不是闭集, 则存在 $E$ 的聚点 $x_0\notin E$. 对于 $\forall x\in E$, 有 $x\ne x_0$, 令
\[
  r_x=\frac12|x-x_0|>0,
\]
则
\[
  U(x,r_x)\cap U(x_0,r_x)=\varnothing.
\]
由于 $\bigcup\limits_{x\in E}U(x,r_x)$ 是 $E$ 的一个开覆盖, 故存在有限子覆盖
\[
  U(x_1,r_{x_1}),U(x_2,r_{x_2}),\cdots,U(x_K,r_{x_K}).
\]
令
\[
  r_K=\min_{1\le k\le K}r_{x_k}>0,
\]
则
\[
  U(x_0,r_K)\cap\left\{\bigcup_{k=1}^K U(x_k,r_{x_k})\right\}=\varnothing.
\]
于是 $U(x_0,r_K)\cap E=\varnothing$, 这与 $x_0$ 是 $E$ 的聚点矛盾. 因此 $E$ 是闭集。

充分性. 设 $E$ 是有界闭集但不是紧集, 则存在 $E$ 的一个开覆盖 $\{O_\lambda\}_{\lambda\in\Lambda}$, 使得它不存在有限子覆盖. 为简单起见, 我们只证明 $E\subset\mathbb R^2$ 的情形, $\mathbb R^n$ 中的情形类似. 由于 $E$ 有界, 存在 $a>0$, 使得 $E\subset[-a,a]\times[-a,a]=I_0$. 将 $I_0$ 等分为四个闭正方形, 则其中至少有一个闭正方形与 $E$ 的交不能被 $\{O_\lambda\}_{\lambda\in\Lambda}$ 中的有限个开集覆盖, 记这个闭正方形为 $I_1$. 再将 $I_1$ 等分为四个闭正方形, 继续上述过程, 得到一列闭正方形 $I_k(k=1,2,\cdots)$, 满足:

(1) $I_k\supset I_{k+1}$;

(2) $\lim\limits_{k\to\infty}\operatorname{diam}(I_k)=0$;

(3) 对于每个 $k$, $I_k\cap E$ 都不能被 $\{O_\lambda\}_{\lambda\in\Lambda}$ 中的有限个开集覆盖。

令 $E_k=I_k\cap E$. 则 $\{E_k\}$ 是满足定理 13.1.10 条件的一列非空闭集. 于是存在唯一的点 $x_0\in\bigcap\limits_{k=1}^{+\infty}E_k$. 因为 $x_0\in E$, 所以存在 $O_{\lambda_0}$, 使得 $x_0\in O_{\lambda_0}$. 又由于 $O_{\lambda_0}$ 是开集, 当 $k$ 充分大时有 $E_k\subset O_{\lambda_0}$, 这与 (3) 矛盾. 证毕。

\section{多元函数与向量函数的极限}

\subsection{多元函数的概念}

在前面的学习中, 我们已经接触到许多一个量依赖于多个变量的情形. 例如, 空间中某点的温度不仅依赖于该点的位置, 还依赖于时间; 火箭在某一时刻达到的高度依赖于推力、重力、燃料质量和火箭自身质量等多个因素. 有了前面的准备, 现在我们可以很自然地引入 $n$ 元函数的概念。

\noindent\textbf{定义 13.2.1}\quad 设 $E\subset\mathbb R^n$. 如果对应法则 $f$ 对于 $E$ 中的每一个点 $x=(x_1,x_2,\cdots,x_n)$, 都有唯一的实数 $u$ 与之对应, 则称 $f$ 是定义在 $E$ 上的 $n$ 元函数, 记为
\[
  f:E\to\mathbb R,\qquad x\mapsto u=f(x)=f(x_1,x_2,\cdots,x_n).
\]
称 $u$ 为函数 $f$ 在点 $x$ 处的函数值, 称 $E$ 为 $f$ 的定义域, 称集合 $\{f(x):x\in E\}\subset\mathbb R$ 为 $f$ 的值域。

当 $n\ge2$ 时, $n$ 元函数也称为多元函数. 一般地, 我们用 $u=f(x)$ 表示多元函数; 必要时也用 $g(x),h(x)$ 等表示. 当 $n=1$ 时, 这与一元函数的记号是一致的, 常写作 $y=f(x)$. 当 $n=2$ 时, 常写作 $z=f(x,y)$; 当 $n=3$ 时, 常写作 $u=f(x,y,z)$。

对于二元函数 $z=f(x,y)$, 如果它具有较好的性质, 那么它的图形是空间中的一个曲面
\[
  \{(x,y,f(x,y)):(x,y)\in E\}.
\]
参见图 13.2.1。

\noindent\textbf{例 13.2.1}\quad 单位球面在 $\mathbb R^3$ 中的方程为
\[
  x^2+y^2+z^2=1.
\]
于是
\[
  z=\sqrt{1-x^2-y^2}
\]
是定义在闭单位圆盘
\[
  \overline\Delta=\{(x,y):x^2+y^2\le1\}
\]
上的二元函数, 它的图形是单位球面的上半球面. 类似地,
\[
  z=-\sqrt{1-x^2-y^2}
\]
的图形是单位球面的下半球面。

设 $u=f(x)(x\in E\subset\mathbb R^n)$. 若 $f(E)$ 是 $\mathbb R$ 中的有界集, 则称 $f$ 在 $E$ 上有界. 类似地, 可以定义函数在 $E$ 中的上界、下界以及上、下确界等. 请读者对一个二元函数来描述上述概念的几何意义。

\subsection{多元函数的极限}

\noindent\textbf{定义 13.2.2}\quad 设函数 $f(x)$ 定义在 $E\subset\mathbb R^n$ 上, $x_0=(x_1^0,x_2^0,\cdots,x_n^0)\in E'$. 若存在 $A\in\mathbb R$, 使得对于 $\forall\varepsilon>0$, 存在 $\delta>0$, 当 $x=(x_1,x_2,\cdots,x_n)\in U_0(x_0,\delta)\cap E$, 即 $0<|x-x_0|<\delta$ 且 $x\in E$ 时, 有
\[
  |f(x)-A|<\varepsilon,
\]
则称 $A$ 为 $x$ 在 $E$ 中趋于 $x_0$ 时函数 $f(x)$ 的极限, 记做
\[
  \lim_{E\ni x\to x_0}f(x)=A,
\]
或
\[
  \lim_{E\ni(x_1,x_2,\cdots,x_n)\to(x_1^0,x_2^0,\cdots,x_n^0)}
  f(x_1,x_2,\cdots,x_n)=A.
\]
如果存在 $x_0$ 的某个去心邻域 $U_0(x_0,\delta_0)\subset E$, 则在极限记号中可以省略 $x\in E$。

上述定义对于一元函数同样适用. 例如, 若 $y=f(x)$ 定义在 $[a,b]$ 上, 则它实际上给出了自变量 $x$ 趋于 $x_0\in(a,b)$ 时的极限以及趋于 $a$ 或 $b$ 时的单侧极限的统一定义. 不仅如此, 上述极限也不要求函数在一个区间上处处有定义, 它的定义域可以是很广泛的集合. 对于多元函数, 上述定义则更为广泛, 因为这时函数的定义域有可能更加复杂。

多元函数极限同一元函数极限一样, 具有唯一性、局部有界性、四则运算法则等性质, 也可以类似地定义无穷小量和无穷大量. 这里不再一一列出。

\noindent\textbf{定理 13.2.1}\quad 设函数 $f(x)$ 定义在 $E\subset\mathbb R^n$ 上, $x_0\in E'$, 则
\[
  \lim_{E\ni x\to x_0}f(x)=A
\]
($A$ 可以为 $\infty$) 的充分必要条件是: 对于 $E\setminus\{x_0\}$ 中任意收敛于 $x_0$ 的点列 $\{x_k\}$, 都有
\[
  \lim_{k\to\infty}f(x_k)=A.
\]

由函数极限的定义知, 当一个函数 $f(x)$ 在 $x_0\in E'$ 存在极限时, 由于 $\mathbb R^n$ 中点 $x$ 趋向 $x_0$ 的方式非常复杂, 因此该极限过程蕴含着许多信息. 反之, 若 $x$ 以某种特殊的方式(如沿曲线、点列等)趋于 $x_0$ 时, $f(x)$ 的极限不存在, 则极限 $\lim\limits_{E\ni x\to x_0}f(x)$ 必不存在。

\noindent\textbf{例 13.2.2}\quad 求极限
\[
  \lim_{(x,y)\to(0,0)}\frac{\tan(x^3+y^3)}{x^2+y^2}.
\]

\noindent\textbf{解}\quad 由于
\[
  \tan(x^3+y^3)\sim x^3+y^3,\qquad (x,y)\to(0,0),
\]
而
\[
  |x^3+y^3|=|x+y||x^2+y^2-xy|
  \le \frac32 |x+y|(x^2+y^2),
\]
故当 $(x,y)\ne(0,0)$ 时, 有
\[
  \left|\frac{x^3+y^3}{x^2+y^2}\right|\le \frac32|x+y|.
\]
由于
\[
  \lim_{(x,y)\to(0,0)}\frac32|x+y|=0,
\]
从而有
\[
  \lim_{(x,y)\to(0,0)}
  \left|\frac{\tan(x^3+y^3)}{x^2+y^2}\right|
  =\lim_{(x,y)\to(0,0)}
  \left|\frac{x^3+y^3}{x^2+y^2}\right|=0.
\]
从上式我们推出
\[
  \lim_{(x,y)\to(0,0)}\frac{\tan(x^3+y^3)}{x^2+y^2}=0.
\]

\noindent\textbf{例 13.2.3}\quad 设函数 $f(x,y)=\dfrac{xy}{x^2+y^2}$, 试问: 当 $(x,y)\to(0,0)$ 时, $f(x,y)$ 是否存在极限?

\noindent\textbf{解法 1}\quad 取点列 $\{(0,\frac1k)\}$, 则有
\[
  \lim_{k\to\infty}f\left(0,\frac1k\right)=\lim_{k\to\infty}0=0.
\]
再取点列 $\{(\frac1k,\frac1k)\}$, 我们有
\[
  \lim_{k\to\infty}f\left(\frac1k,\frac1k\right)
  =\lim_{k\to\infty}\frac{\frac1{k^2}}{2\frac1{k^2}}=\frac12.
\]
因此当 $(x,y)$ 趋于 $(0,0)$ 时, $f(x,y)$ 不存在极限。

\noindent\textbf{解法 2}\quad 取 $y=kx(k\in\mathbb R)$, 则当 $x\to0$ 时, 有 $y\to0$. 由于
\[
  \lim_{x\to0}f(x,kx)
  =\lim_{x\to0}\frac{kx^2}{(1+k^2)x^2}
  =\frac{k}{1+k^2},
\]
这说明自变量 $(x,y)$ 沿射线 $y=kx$ 趋于 $(0,0)$ 时, $f(x,y)$ 的极限依赖于射线的选取, 从而极限 $\lim\limits_{(x,y)\to(0,0)}f(x,y)$ 必不存在。

\noindent\textbf{解法 3}\quad 在 $Oxy$ 平面上取极坐标
\[
  \begin{cases}
    x=r\cos\theta,\\
    y=r\sin\theta,
  \end{cases}
\]
其中 $r\in(0,+\infty),\theta\in[0,2\pi)$, 则 $f(x,y)$ 在极坐标下为
\[
  f(x,y)=\frac12\sin2\theta.
\]
由于 $(x,y)\to(0,0)$ 时必有 $r\to0+0$, 而 $\frac12\sin2\theta$ 与 $r$ 无关但又不是常数函数, 从而极限 $\lim\limits_{(x,y)\to(0,0)}f(x,y)$ 也不存在。

\noindent\textbf{例 13.2.4}\quad 设函数
\[
  f(x,y)=\frac{x^2y}{x^4+y^2},
\]
讨论 $\lim\limits_{(x,y)\to(0,0)}f(x,y)$ 是否存在。

\noindent\textbf{解}\quad 如果取 $y=kx(k\in\mathbb R)$, 则当 $(x,y)$ 沿该射线趋于 $(0,0)$ 时, 有
\[
  \lim_{x\to0}f(x,kx)=\lim_{x\to0}\frac{kx^3}{(x^2+k^2)x^2}=0;
\]
如果取 $y=x^2$, 则当 $(x,y)$ 沿该抛物线趋于 $(0,0)$ 时, 有
\[
  \lim_{x\to0}f(x,x^2)=\lim_{x\to0}\frac{x^4}{x^4+x^4}=\frac12.
\]
这说明 $\lim\limits_{(x,y)\to(0,0)}f(x,y)$ 不存在。

此例说明, 一个函数沿任何射线的极限都存在且相等也不足以保证该函数的极限存在。

\noindent\textbf{例 13.2.5}\quad 设函数
\[
  f(x)=f(x,y,z)=[\sin1-\sin(x^2+y^2+z^2)]\ln[1-(x^2+y^2+z^2)].
\]
记 $E=\{x=(x,y,z):|x|<1\}$, 并设 $x_0\in\partial E$, 求极限 $\lim\limits_{E\ni x\to x_0}f(x)$。

\noindent\textbf{解}\quad 显然 $f$ 的定义域为 $E$. 当 $E\ni x\to x_0$ 时, $t=x^2+y^2+z^2\to1-0$, 因此
\[
  \lim_{E\ni x\to x_0}f(x)
  =\lim_{t\to1-0}(\sin1-\sin t)\ln(1-t)=0.
\]
最后一步的极限可由求一元函数极限的方法求得(如用洛必达法则)。

\subsection{累次极限}

设 $f(x)=f(x_1,x_2,\cdots,x_n)$ 在 $x_0\in\mathbb R^n$ 的邻域内有定义, 现在继续考虑 $x$ 趋于 $x_0$ 时极限行为的方式. 当 $x$ 趋于 $x_0$ 时具有许多种方式, 且 $x$ 沿着某些特殊方式趋于 $x_0$ 时具有特别的意义. 在以后的章节中我们常常会遇到以下形式的累次极限. 我们以一个二元函数的情形为例来讲述这种极限, 对 $n(n\ge3)$ 元函数的情形是类似的, 只是形式复杂而已。

\noindent\textbf{定义 13.2.3}\quad 设函数 $z=f(x,y)$ 在 $E\subset\mathbb R^2$ 上有定义, 且邻域 $N_0((x_0,y_0),\delta_0)\subset E(\delta_0>0)$. 若在 $N_0((x_0,y_0),\delta_0)$ 内, 对每个固定的 $y\ne y_0$, $\lim\limits_{x\to x_0}f(x,y)=\varphi(y)$ 存在, 并且 $\lim\limits_{y\to y_0}\varphi(y)=A(A$ 为常数), 则称 $A$ 为 $f(x,y)$ 当 $(x,y)$ 趋于 $(x_0,y_0)$ 时先 $x$ 后 $y$ 的累次极限, 记为
\[
  A=\lim_{y\to y_0}\lim_{x\to x_0}f(x,y).
\]
类似地, 我们可以定义先 $y$ 后 $x$ 的累次极限 $\lim\limits_{x\to x_0}\lim\limits_{y\to y_0}f(x,y)$。

从定义 13.2.3 可以看出, 一个二元函数在 $(x_0,y_0)$ 处的累次极限存在与否与函数在直线 $x=x_0$ 与 $y=y_0$ 上的值无关. 因此容易举出例子来说明一个函数的两个累次极限都存在, 但有可能极限却不存在. 例如, 设函数
\[
  f(x,y)=
  \begin{cases}
    0, & xy=0,\\
    1, & xy\ne0,
  \end{cases}
\]
则有
\[
  \lim_{x\to0}\lim_{y\to0}f(x,y)
  =\lim_{y\to0}\lim_{x\to0}f(x,y)=1,
\]
而 $\lim\limits_{(x,y)\to(0,0)}f(x,y)$ 却不存在。

另外, 当函数的极限存在时, 人们很自然地会猜测累次极限必定存在且等于函数的极限. 对于该猜测, 我们的答案是否定的. 例如, 设
\[
  f(x,y)=
  \begin{cases}
    (x+y)\sin\dfrac1x\sin\dfrac1y, & xy\ne0,\\
    0, & xy=0.
  \end{cases}
\]
由于对于 $\forall(x,y)\in\mathbb R^2$, 有 $|f(x,y)|\le |x|+|y|$, 因此
\[
  \lim_{(x,y)\to(0,0)}f(x,y)=0.
\]
注意到对每个 $x\ne0,\frac1{k\pi}(k$ 为整数), 极限 $\lim\limits_{y\to0}f(x,y)$ 不存在; 而对每个 $y\ne0,\frac1{k\pi}(k$ 为整数), 极限 $\lim\limits_{x\to0}f(x,y)$ 也不存在. 因此无法讨论 $(x,y)$ 趋于 $(0,0)$ 时的累次极限, 所以 $f(x,y)$ 的两个累次极限都不存在。

考查以上例子, 其累次极限不存在的本质是存在任意小的 $x\ne0$, 极限 $\lim\limits_{y\to0}f(x,y)$ 不存在. 若假定后者处处存在, 则可以证明人们的猜测是正确的。

\noindent\textbf{定理 13.2.2}\quad 设函数 $f(x,y)$ 在 $N_0((x_0,y_0),\delta_0)(\delta_0>0)$ 上有定义, 并且满足:

(1) $\exists A$(可以是 $\infty$), 使得
\[
  \lim_{(x,y)\to(x_0,y_0)}f(x,y)=A;
\]

(2) 对于 $U_0(y_0,\delta_0)$ 内任意固定的 $y\ne y_0$, 极限 $\lim\limits_{x\to x_0}f(x,y)=\varphi(y)$ 存在,

则有
\[
  \lim_{y\to y_0}\varphi(y)
  =\lim_{y\to y_0}\lim_{x\to x_0}f(x,y)=A.
\]

\noindent\textbf{证明}\quad 我们只证 $A\ne\infty$ 的情形。

由 $\lim\limits_{(x,y)\to(x_0,y_0)}f(x,y)=A$ 可知, 对于 $\forall\varepsilon>0,\exists\delta(0<\delta<\delta_0)$, 当 $(x,y)\in N_0((x_0,y_0),\delta)$ 时, 有
\begin{equation}
  |f(x,y)-A|<\frac{\varepsilon}{2}.
  \tag{13.2.1}
\end{equation}
对任意的 $y\ne y_0$, 当 $(x,y)\in N((x_0,y_0),\delta)$ 时, 由于 $\lim\limits_{x\to x_0}f(x,y)=\varphi(y)$ 存在, 因此可在式 (13.2.1) 中令 $x\to x_0$, 从而当 $0<|y-y_0|<\delta$ 时, 有
\[
  |\varphi(y)-A|\le\frac{\varepsilon}{2}<\varepsilon,
\]
即 $\lim\limits_{y\to y_0}\varphi(y)=A$. 证毕。

\subsection{向量函数的定义与极限}

设 $A,B$ 为任意两个非空集合(不一定是数集), 若对于 $A$ 中每个元素 $a$, 根据对应法则 $f$, 在 $B$ 中有唯一一个元素 $b$ 与之对应, 则称 $f$ 是 $A$ 到 $B$ 的一个映射, $b$ 称为 $a$ 在映射 $f$ 下的像, 记为
\[
  f:A\to B,\qquad a\mapsto b=f(a).
\]
类似于函数的有关定义, 我们可以给出一个映射的定义域、值域等概念. 同样, 对于映射我们可以引进单射、满射以及一一对应等概念. 利用映射的概念, 容易给出向量函数的定义。

\noindent\textbf{定义 13.2.4}\quad 设 $E\subset\mathbb R^n$, 若 $f$ 是 $E$ 到 $\mathbb R^m(m\ge1)$ 的一个映射, 则称 $f$ 是 $E$ 上的一个 $m$ 维 $(n$ 元)向量函数(简称向量函数), 记为
\[
  u=f(x),\qquad x\in E.
\]

设 $E\subset\mathbb R^n$, 显然, 一个 $E\to\mathbb R$ 的向量函数即是前面我们讨论的 $n$ 元函数. 对于上述定义的 $m$ 维向量函数, 记 $x=(x_1,x_2,\cdots,x_n)\in E$, $u=(u_1,u_2,\cdots,u_m)=f(x_1,x_2,\cdots,x_n)\in\mathbb R^m$, 则 $f(x)$ 的每个分量 $u_j$ 都是 $E$ 上的一个 $n$ 元函数, 即对于 $j=1,2,\cdots,m$, 有 $u_j=f_j(x)(x\in E)$. 因此 $u=f(x):E\to\mathbb R^m$ 为一个向量函数的充分必要条件是: 存在 $m$ 个定义在 $E$ 上的 $n$ 元函数 $f_j(x)(j=1,2,\cdots,m)$, 使得
\[
  f(x)=(f_1(x),f_2(x),\cdots,f_m(x)).
\]
自然地, 若 $f(E)$ 是 $\mathbb R^m$ 的有界集, 则称 $f(x)$ 在 $E$ 是有界的。

例如, 平面 $\mathbb R^2$ 中子集
\[
  E=\{(r,\theta):0\le r<+\infty,\ 0\le\theta<2\pi\}
\]
到 $\mathbb R^2$ 上的极坐标变换 $f(r,\theta)=(x,y)=(r\cos\theta,r\sin\theta)$ 是一个二维(二元)向量函数. 作为映射, 容易看出它是 $E\setminus\{(0,\theta):0\le\theta<2\pi\}$ 到 $\mathbb R^2\setminus\{(0,0)\}$ 的一一对应. 在 $(0,0)$ 处它不是单射. 事实上, 对于 $\forall\theta(0\le\theta<2\pi)$, 我们有 $f(0,\theta)=(0,0)$。

\noindent\textbf{定义 13.2.5}\quad 设 $E\subset\mathbb R^n,x_0\in E'$, $f(x)=(f_1(x),f_2(x),\cdots,f_m(x))$ 是定义在 $E$ 上的一个 $m$ 维向量函数. 若对于 $\forall j(1\le j\le m)$, 有
\[
  \lim_{E\ni x\to x_0}f_j(x)=A_j,
\]
其中 $A_j\in\mathbb R(j=1,2,\cdots,m)$, 则称 $f(x)$ 当 $x$(在 $E$ 中)趋于 $x_0$ 时的极限为 $(A_1,A_2,\cdots,A_m)$, 记为
\[
  \lim_{E\ni x\to x_0}f(x)=(A_1,A_2,\cdots,A_m).
\]
若 $\exists j_0(1\le j_0\le m)$, $\lim\limits_{E\ni x\to x_0}f_{j_0}(x)$ 不存在, 则称 $f(x)$ 当 $x$(在 $E$ 中)趋于 $x_0$ 时极限不存在。

\noindent\textbf{例 13.2.6}\quad 设向量函数 $f(r,\theta)=(r\cos\theta,r\sin\theta)$, 其中 $r\in(0,+\infty),\theta\in[0,2\pi)$, 再设 $r_0>0$, 求极限
\[
  \lim_{(r,\theta)\to(r_0,0)}f(r,\theta)
  \quad\text{和}\quad
  \lim_{(r,\theta)\to(r_0,2\pi)}f(r,\theta).
\]

\noindent\textbf{解}\quad 由定义 13.2.5 我们有
\[
  \lim_{(r,\theta)\to(r_0,0)}f(r,\theta)
  =\lim_{(r,\theta)\to(r_0,0)}(r\cos\theta,r\sin\theta)=(r_0,0)
\]
和
\[
  \lim_{(r,\theta)\to(r_0,2\pi)}f(r,\theta)
  =\lim_{(r,\theta)\to(r_0,2\pi)}(r\cos\theta,r\sin\theta)=(r_0,0).
\]

\section{多元连续函数}

\subsection{多元连续函数}

与一元函数类似, 我们可以引进多元函数连续性的相关概念。

\noindent\textbf{定义 13.3.1}\quad 设函数 $u=f(x)$ 在 $E\subset\mathbb R^n$ 上有定义, $x_0\in E$. 若 $x_0$ 是 $E$ 的孤立点或当 $x_0\in E'$ 时有
\[
  \lim_{E\ni x\to x_0}f(x)=f(x_0),
\]
则称 $f(x)$ 在 $x_0$ 处连续. 若 $f(x)$ 在 $E$ 中的每一点都连续, 则称 $f(x)$ 在 $E$ 上连续. 若 $\lim\limits_{E\ni x\to x_0}f(x)$ 不存在, 或存在但不等于 $f(x_0)$, 则称 $f(x)$ 在 $x_0$ 处不连续或间断, 此时 $x_0$ 也称做 $f(x)$ 的间断点。

首先, 当 $x_0$ 是 $E$ 的孤立点时, 我们自然地认为 $f(x)$ 在 $x_0$ 处连续. 当 $x_0$ 是 $E$ 的内点时, $f(x)$ 在 $x_0$ 处的连续性是我们在一元函数里通常见到的情形, 但是以上定义还包括了非常广泛的情形, 例如一元函数在闭区间的连续的定义等. 读者可以容易将一元连续函数的一些局部性质(即局部有界、保号等)以及和、差、积、商等结果平行推广到多元连续函数的情形。

\noindent\textbf{例 13.3.1}\quad 讨论函数
\[
  f(x,y)=
  \begin{cases}
    x\sin\dfrac1y, & xy\ne0,\\
    0, & xy=0
  \end{cases}
\]
在 $\mathbb R^2$ 中的连续性。

\noindent\textbf{解}\quad 显然当 $(x_0,y_0)\in\mathbb R^2$ 且 $x_0y_0\ne0$ 时, $f(x,y)$ 在 $(x_0,y_0)$ 处连续. 由于对于 $\forall y\in(-\infty,+\infty)\setminus\{0\}$, 有
\[
  \left|x\sin\frac1y\right|\le |x|,
\]
注意到当 $y=0$ 时, 有 $f(x,0)=0$, 于是对于 $\forall y_0\in(-\infty,+\infty)$, 有
\[
  \lim_{(x,y)\to(0,y_0)}f(x,y)=0=f(0,y_0).
\]
这说明 $f(x,y)$ 在 $y$ 轴上连续。

当 $x_0\ne0$ 时, 由于
\[
  \lim_{(x_0,y)\to(x_0,0)}f(x_0,y)
  =\lim_{(x_0,y)\to(x_0,0)}x_0\sin\frac1y
\]
不存在, 所以 $\lim\limits_{(x,y)\to(x_0,0)}f(x,y)$ 也不存在. 因此 $f(x,y)$ 在 $\mathbb R^2$ 中的间断点集是正负实轴构成的集合, 即 $\{(x,y):x\ne0,y=0\}$。

\subsection{多元连续向量函数}

对于多元向量函数, 我们自然地引进下述定义。

\noindent\textbf{定义 13.3.2}\quad 设 $f(x)=(f_1(x),f_2(x),\cdots,f_m(x))$ 是 $E\subset\mathbb R^n$ 到 $\mathbb R^m$ 的一个 $m$ 维向量函数, 并设 $x_0\in E$. 若 $x_0$ 是 $E$ 的孤立点或当 $x_0\in E'$ 时有
\[
\begin{aligned}
  \lim_{E\ni x\to x_0}f(x)
  &=
  \left(\lim_{E\ni x\to x_0}f_1(x),
  \lim_{E\ni x\to x_0}f_2(x),
  \cdots,\lim_{E\ni x\to x_0}f_m(x)\right)\\
  &=(f_1(x_0),f_2(x_0),\cdots,f_m(x_0))=f(x_0),
\end{aligned}
\]
则称 $f(x)$ 在 $x_0$ 处连续. 若 $f(x)$ 在 $E$ 中的每一点都连续, 则称 $f(x)$ 在 $E$ 上连续. 若存在 $j_0(1\le j_0\le m)$, $f_{j_0}(x)$ 在 $x_0$ 处不连续, 则称 $f(x)$ 在 $x_0$ 处不连续, 并称 $x_0$ 是 $f(x)$ 的间断点。

由定义可知, $f(x)$ 在 $x_0$ 处连续的充分必要条件是它的每个分量函数 $f_j(x)(j=1,2,\cdots,m)$ 都在 $x_0$ 处连续. 容易验证, 对于定义域相同的两个多元连续向量函数, 它们的和、差以及数乘还是多元连续向量函数。

下面我们不加证明地给出经常用到的多元复合向量函数的连续性定理。

\noindent\textbf{定理 13.3.1}\quad 设 $E\subset\mathbb R^n$, 向量函数 $y=f(x)=(f_1(x),f_2(x),\cdots,f_m(x))$ 在 $E$ 上连续; $g(y)=(g_1(y),g_2(y),\cdots,g_l(y))$ 在 $f(E)\subset\mathbb R^m$ 上连续, 则
\[
\begin{aligned}
  g(f(x))={}&(g_1(f_1(x),f_2(x),\cdots,f_m(x)),\\
  &g_2(f_1(x),f_2(x),\cdots,f_m(x)),\\
  &\cdots,\ g_l(f_1(x),f_2(x),\cdots,f_m(x)))
\end{aligned}
\]
在 $E$ 上连续。

\noindent\textbf{例 13.3.2}\quad 设函数
\[
  f(x,y)=
  \begin{cases}
    \sqrt{1-x^2-y^2}\ln(1-x^2-y^2), & x^2+y^2<1,\\
    0, & x^2+y^2=1,
  \end{cases}
\]
讨论 $f(x,y)$ 的连续性。

\noindent\textbf{解}\quad 记 $\Delta=\{(x,y):x^2+y^2<1\}$, 则 $f(x,y)$ 的定义域是 $\overline\Delta$. 任取 $(x_0,y_0)\in\Delta$, 容易看出 $1-x^2-y^2$ 在 $(x_0,y_0)$ 处连续且值大于零, 从而由复合函数的连续性知 $\sqrt{1-x^2-y^2}$ 及 $\ln(1-x^2-y^2)$ 在 $(x_0,y_0)$ 处连续. 所以 $f(x,y)$ 在 $(x_0,y_0)$ 处连续. 当 $(x_0,y_0)\in\partial\Delta$ 且 $\Delta\ni(x,y)\to(x_0,y_0)$ 时, 我们有 $t=x^2+y^2\to1-0$. 由于
\[
  \lim_{t\to1-0}\sqrt{1-t}\ln(1-t)=0,
\]
所以
\[
  \lim_{\Delta\ni(x,y)\to(x_0,y_0)}
  \sqrt{1-x^2-y^2}\ln(1-x^2-y^2)=0=f(x_0,y_0).
\]
另外, 我们显然有
\[
  \lim_{\partial\Delta\ni(x,y)\to(x_0,y_0)}f(x,y)=0=f(x_0,y_0),
\]
从而有
\[
  \lim_{\overline\Delta\ni(x,y)\to(x_0,y_0)}f(x,y)=0=f(x_0,y_0).
\]
这说明 $f(x,y)$ 在 $\overline\Delta$ 上连续。

\subsection{集合的连通性}

设 $D=[\alpha,\beta]\subset\mathbb R(\alpha<\beta)$, 我们通常称 $D$ 上的一个 $n$ 维连续函数
\[
  h(t)=(x_1(t),x_2(t),\cdots,x_n(t))
\]
为 $\mathbb R^n$ 中的一条连续曲线(简称曲线). 当 $n=1,2,3$ 时, 连续曲线具有鲜明的几何意义. 一条连续曲线也称为 $\mathbb R^n$ 中的一条道路(或路径)。

\noindent\textbf{定义 13.3.3}\quad 设 $E\subset\mathbb R^n$ 是一个非空集合, 若对于 $\forall x,y\in E$, 都存在一条道路 $h(t):t\in[\alpha,\beta]$, 使得 $h(\alpha)=x,h(\beta)=y$, 且对于 $\forall t\in[\alpha,\beta]$, 有 $h(t)\in E$, 则称 $E$ 是道路连通的。

\noindent\textbf{注}\quad 对于一般集合的连通性在后续课程中有更为广泛的定义与讨论. 但在本书中, 我们提到的连通性则是指上述定义中的道路连通. 今后, 我们将省略``道路''两字, 仅仅说一个集合是否是连通的。

在 $\mathbb R$ 中, 一个集合 $E$ 连通的充分必要条件是 $E$ 是一个区间. 在 $\mathbb R^2$ 或 $\mathbb R^3$ 中, 若一个集合 $E$ 是连通的, 则它的形式可以多种多样. 特别地, 若 $E$ 是连通的开集, 则直观地看 $E$ 只能由``一块''组成。

若 $E\subset\mathbb R^n$ 是一个连通的开集, 则称 $E$ 为一个区域. 一般用 $D,\Omega$ 等来表示区域. 设 $D$ 是一个区域, 称 $D\cup\partial D$ 是一个闭区域. 今后我们也经常用 $D,\Omega$ 等来表示一个闭区域. 例如, 以原点为心, 1 为半径的单位圆盘 $\Delta=\{(x,y):x^2+y^2<1\}$ 与 $\mathbb R^2\setminus\overline\Delta$ 都是区域. 在今后的学习中, 我们遇到比较多的是 $\mathbb R^2$ 和 $\mathbb R^3$ 中的区域. 一般来说, 在微积分课程中, 我们在 $\mathbb R^2$ 中遇到的有界区域一般是由有限多条分段光滑曲线所围的区域. 如图 13.3.1 所示的区域 $D$, 它是由 $K+1$ 条封闭曲线所围成的有界区域. 而在 $\mathbb R^3$ 中, 我们则主要考虑由若干块曲面所围的区域. 例如, 单位球面 $x^2+y^2+z^2=1$ 的内部位于 $Oxy$ 平面上方的部分就是一个区域, 它由单位球面的上半部分与 $Oxy$ 平面所围成。

区域在多元微积分中所起的作用与开区间在一元微积分中所起的作用相似, 而闭区域在多元微积分中所起的作用与闭区间在一元微积分中所起的作用相似。

今后我们还会遇到凸域的概念. 设 $D\subset\mathbb R^n$ 是一个区域, 若对于 $\forall x,y\in D$, 有 $tx+(1-t)y\in D(t\in[0,1])$, 即连接 $x$ 与 $y$ 的线段在 $D$ 内, 则称 $D$ 为凸域。

\subsection{连续函数的性质}

在本小节中, 我们将闭区间上一元连续函数的性质推广到多元(或向量)函数的情形. 将一个已知数学定理作推广, 考查该定理的证明过程非常重要. 例如, 对于有界闭区间上的一元连续函数必有界并取到最大、最小值这一定理, 如果利用有界序列存在收敛子列这一结论来证的话, 则只要求该闭区间是紧集即可. 而用类似的论证方法很容易就可以证明高维情形的相应结果。

\noindent\textbf{定理 13.3.2}\quad 设 $E\subset\mathbb R^n$ 为一个紧集, 向量函数
\[
  u=f(x)=(f_1(x),f_2(x),\cdots,f_m(x))
\]
在 $E$ 上连续, 则 $f(E)$ 是 $\mathbb R^m$ 中的紧集。

\noindent\textbf{证明}\quad 首先我们证明 $f(E)$ 是有界集合. 倘若结论不真, 则存在 $E$ 中的点列 $\{x_k\}$, 使得
\[
  \lim_{k\to\infty}|f(x_k)|=+\infty.
\]
由 $E$ 的紧性, 存在 $\{x_k\}$ 的子序列 $\{x_{k_j}\}$, 使得
\[
  \lim_{j\to\infty}x_{k_j}=x_0\in E.
\]
再由 $f(x)$ 在 $E$ 上的连续性知
\[
  \lim_{j\to\infty}f(x_{k_j})=f(x_0).
\]
此矛盾证明了 $f(E)$ 的有界性。

再证明 $f(E)$ 是闭集. 设 $u_0$ 是 $f(E)$ 的一个聚点, 我们要证 $u_0\in f(E)$. 由于 $u_0$ 是 $f(E)$ 的聚点, 从而存在 $\{x_k'\}\subset E$, 使得
\[
  \lim_{k\to\infty}f(x_k')=u_0.
\]
再由 $E$ 是紧集知, $\{x_k'\}$ 必存在收敛子列 $\{x_{k_j}'\}$. 设
\[
  \lim_{j\to\infty}x_{k_j}'=x_0',
\]
则有 $x_0'\in E$. 最后由 $f(x)$ 在 $x_0'$ 处连续, 有
\[
  u_0=\lim_{j\to\infty}f(x_{k_j}')=f(x_0')\in f(E).
\]
证毕。

由于对 $\mathbb R$ 中的紧集 $E$ 必有 $\sup E=\max E$ 和 $\inf E=\min E$, 从定理 13.3.2 我们可以推出下面的结论。

\noindent\textbf{推论}\quad 设 $E\subset\mathbb R^n$ 为一个紧集, 函数 $f(x)$ 在 $E$ 上连续, 则

(1) $f(x)$ 在 $E$ 上有界;

(2) $f(x)$ 在 $E$ 上取到最大、最小值。

一个区间上的一元连续函数具有介值性质. 这一性质成立的充分必要条件是该函数的值域是一个区间, 而这正是连通集合的特征. 对于多元向量函数, 我们有以下定理。

\noindent\textbf{定理 13.3.3}\quad 设 $E\subset\mathbb R^n$ 是连通集, 向量函数 $f(x)=(f_1(x),f_2(x),\cdots,f_m(x))$ 在 $E$ 上连续, 则 $f(E)$ 是 $\mathbb R^m$ 中的连通集。

\noindent\textbf{证明}\quad 设 $u_1,u_2\in f(E)\subset\mathbb R^m$, 则 $\exists x_1,x_2\in E$, 使得
\[
  u_j=f(x_j)\qquad (j=1,2).
\]
由 $E$ 是连通的, 存在连接 $x_1,x_2$ 的道路 $h(t)=(x_1(t),x_2(t),\cdots,x_n(t))(t\in[0,1])$, 即对于 $\forall t\in[0,1]$, 有 $h(t)\in E$, 且 $h(0)=x_1,h(1)=x_2$. 容易看出
\[
\begin{aligned}
  f(h(t))={}&(f_1(x_1(t),x_2(t),\cdots,x_n(t)),\\
  &f_2(x_1(t),x_2(t),\cdots,x_n(t)),\\
  &\cdots,\ f_m(x_1(t),x_2(t),\cdots,x_n(t)))
\end{aligned}
\]
是 $f(E)$ 中连接 $u_1,u_2$ 的一条道路. 这说明 $f(E)$ 是 $\mathbb R^m$ 中的连通集. 证毕。

\noindent\textbf{推论}\quad 设 $E\subset\mathbb R^n$ 是连通集, 函数 $u=f(x)$ 在 $E$ 上连续, 再设 $u_1,u_2\in f(E)$, 且 $u_1<u_2$, 则对于 $\forall c\in(u_1,u_2)$, 存在 $\xi\in E$, 使得
\[
  f(\xi)=c.
\]

\noindent\textbf{证明}\quad 取 $x_1,x_2\in E$, 使得 $u_1=f(x_1),u_2=f(x_2)$, 并任取 $E$ 中的道路 $h(t)(0\le t\le1)$ 连接 $x_1,x_2$, 则 $f(h(t))$ 是 $[0,1]$ 上的连续函数. 由定理 13.3.1(或一元连续函数的介值定理)知, 存在 $\eta\in(0,1)$, 使得
\[
  f(h(\eta))=c.
\]
记 $\xi=h(\eta)\in E$, 即得 $f(\xi)=c$. 证毕。

在上述推论中, 当 $E$ 是 $\mathbb R^n(n\ge2)$ 的一个区域时, 满足 $f(\xi)=c$ 的集合 $\{\xi\in E:f(\xi)=c\}$ 一般有无穷多个点. 这是因为 $E$ 中连接两点可以有无穷多条道路并且任两条道路仅仅只在端点重合。

设 $f(x)(x\in E\subset\mathbb R^n)$ 是一个 $n$ 元函数, 如果对于 $\forall\varepsilon>0,\exists\delta>0$, 当 $x',x''\in E$ 且 $|x'-x''|<\delta$ 时, 有
\[
  |f(x')-f(x'')|<\varepsilon,
\]
则称 $f(x)$ 在 $E$ 上一致连续. 设 $f(x)=(f_1(x),f_2(x),\cdots,f_m(x))(x\in E\subset\mathbb R^n)$ 是一个 $m$ 维向量函数, 如果对于 $\forall j(1\le j\le m)$, $f_j(x)$ 在 $E$ 上一致连续, 则称 $f(x)$ 在 $E$ 上一致连续。

\noindent\textbf{定理 13.3.4}\quad 设 $E\subset\mathbb R^n$ 是紧集, 若向量函数 $f(x)$ 在 $E$ 上连续, 则 $f(x)$ 在 $E$ 一致连续。

定理 13.3.4 的证明与有界闭区间上的一元连续函数的情形类似, 请读者自证。

\noindent\textbf{例 13.3.3}\quad 设函数 $f(x,y)=\sin xy$, 证明: $f(x,y)$ 在 $\mathbb R^2$ 中的任何紧集上是一致连续的, 但它在 $\mathbb R^2$ 中不是一致连续的。

\noindent\textbf{证明}\quad 显然 $f(x,y)=\sin xy$ 在 $\mathbb R^2$ 内处处连续, 由定理 13.3.4 知 $f(x,y)$ 在 $\mathbb R^2$ 中的任何紧集上是一致连续的。

下证 $f(x,y)$ 在 $\mathbb R^2$ 中不一致连续. 事实上, 对 $k=1,2,\cdots$, 令
\[
  x_k'=\left(k,\frac1k\right)\quad\text{和}\quad
  x_k''=\left(k,\frac2k\right),
\]
则
\[
  \lim_{k\to\infty}|x_k'-x_k''|=\lim_{k\to\infty}\frac1k=0.
\]
但对于 $\forall k\in\mathbb N$, 我们有
\[
  \left|f\left(k,\frac1k\right)-f\left(k,\frac2k\right)\right|
  =|\sin1-\sin2|\ne0.
\]
这就证明了 $f(x,y)$ 在 $\mathbb R^2$ 中不一致连续。

\subsection{同胚映射}

设 $E\subset\mathbb R^n$, $y=f(x)=(f_1(x),f_2(x),\cdots,f_m(x))$ 是定义在 $E$ 上的一个向量函数. 作为映射, 若 $f(x):E\to f(E)$ 是一个一一对应, 则存在逆映射
\[
  x=f^{-1}(y):f(E)\to E.
\]
如果 $f(x)$ 在 $E$ 上连续以及 $f^{-1}(y)$ 在 $f(E)$ 上连续, 则称 $y=f(x)$ 是 $E\to f(E)$ 的同胚映射(简称同胚). 同胚映射 $f(x)$ 也称为 $E$ 到 $f(E)$ 的变换, 它的逆映射也称为逆变换。

\noindent\textbf{例 13.3.4}\quad 证明: 不存在 $\mathbb R$ 到 $\mathbb R^2$ 的同胚映射。

\noindent\textbf{证明}\quad 用反证法. 倘若存在 $\mathbb R$ 到 $\mathbb R^2$ 的同胚映射 $y=f(x)$. 取 $E=\mathbb R\setminus\{0\}$, 则 $E$ 是 $\mathbb R$ 中的不连通集. 令 $\mathbb R^2\setminus\{f(0)\}=f(E)$, 则 $f(E)$ 在 $\mathbb R^2$ 中仍然是连通集. 由于 $f^{-1}(y)$ 在 $f(E)$ 上连续, 从而它将 $f(E)$ 映成 $\mathbb R$ 中的连通集, 但 $f^{-1}(y)$ 将 $f(E)$ 映成 $E$. 此矛盾便证明了我们的结论。

对于一个区间上的一元连续函数, 若它存在反函数, 则其反函数必定连续. 这是因为在一元函数的情形, 我们涉及的函数具有单调性的缘故. 对于多元函数, 则上述结论可以不真。

事实上, 考虑极坐标变换
\[
  (x,y)=f(r,\theta)=(r\cos\theta,r\sin\theta),
\]
则 $f(r,\theta)$ 将
\[
  D=\{(r,\theta):0<r<+\infty,\ 0\le\theta<2\pi\}
\]
连续地并且一一地映成 $\Omega=\mathbb R^2\setminus\{(0,0)\}$, 因此它的逆映射 $(r,\theta)=f^{-1}(x,y)$ 存在. 我们可以清楚地写出它的表达式:
\[
  r=\sqrt{x^2+y^2},
\]
\[
  \theta=
  \begin{cases}
    \arctan(y/x), & x>0,\ y\ge0,\\
    \pi/2, & x=0,\ y>0,\\
    \pi+\arctan(y/x), & x<0,\\
    3\pi/2, & x=0,\ y<0,\\
    2\pi+\arctan(y/x), & x>0,\ y<0.
  \end{cases}
\]
从几何上来看, 对于极坐标变换
\[
  f:
  \begin{cases}
    x=r\cos\theta,\\
    y=r\sin\theta
  \end{cases}
\]
来说, 例 13.2.6 告诉我们, 若取定 $r_0>0$, 则当 $(r,\theta)$ 趋于 $(r_0,0)$ 时, 对应的 $(x,y)$ 从第一象限趋于 $(r_0,0)$; 而当 $(r,\theta)$ 趋于 $(r_0,2\pi)$ 时, 对应的 $(x,y)$ 从第四象限趋于 $(r_0,0)$. 这说明当 $(x,y)\to(r_0,0)$ 时, $(r,\theta)=f^{-1}(x,y)$ 不存在极限, 即 $f^{-1}(x,y)$ 在正实轴上处处不连续。

从上述的例子我们可以看出, 极坐标变换的定义域 $D$ 不是区域, 而 $f(D)$ 是一个区域. 若取 $D=\{(r,\theta):0<r<+\infty,\ 0<\theta<2\pi\}$, 则
\[
  f(D)=\mathbb R^2\setminus\{(x,y):x\ge0,\ y=0\},
\]
即 $f(D)$ 为 $\mathbb R^2$ 中挖去正 $x$ 轴与原点的区域. 此时读者不难看出, $f(r,\theta)$ 是 $D$ 到 $f(D)$ 的变换。

\section*{习题十三}

1. 证明 $\mathbb R^n$ 中两点间的距离满足三角不等式: 对于 $\forall x,y$ 和 $z\in\mathbb R^n$, 成立
\[
  |x-z|\le |x-y|+|y-z|.
\]

2. 若 $\lim\limits_{k\to\infty}|x_k|=+\infty$, 则称 $\mathbb R^n$ 中的点列 $\{x_k\}$ 趋于 $\infty$. 现在设点列 $\{x_k=(x_1^k,x_2^k,\cdots,x_n^k)\}$ 趋于 $\infty$, 试判断下列命题是否正确:

(1) 对于 $\forall i(1\le i\le n)$, 序列 $\{x_i^k\}$ 趋于 $\infty$;

(2) $\exists i_0(1\le i_0\le n)$, 序列 $\{x_{i_0}^k\}$ 趋于 $\infty$。

3. 求下列集合的聚点集:
\[
  (1)\quad E=\left\{\left(\frac qp,\frac qp,1\right)\in\mathbb R^3:p,q\in\mathbb N\text{ 互素, 且 }q<p\right\};
\]
\[
  (2)\quad E=\left\{\left(\left(\ln\left(1+\frac1k\right)\right)^k,\sin\frac{k\pi}{2}\right):k=1,2,\cdots\right\}.
\]

\end{document}
